\exercisesection{4}

\begin{enumerate}
\def\labelenumi{\arabic{enumi}.}
\tightlist
\item
  设 A 为交换环，带滤过 \((a_{n})_{n\geq0}\) 且关于此滤过是 Hausdorff 且完备的，满足 \(a_{0}=A\)。令 M、\(M'\)、N 为三个 A-模，每个都赋予由 A 上滤过诱导的滤过以及由该滤过定义的拓扑；记 \(\overline{M}=M/a_{1}M\)、\(\overline{M}'=M'/a_{1}M'\)、\(\overline{N}=N/a_{1}N\)。令 \(f: M\times M'\to N\) 为双线性映射，\(\bar{f}: \overline{M}\times\overline{M'}\to\overline{N}\) 为由 f 取商得到的 \((A/a_{1})\)-双线性映射。令 \(y\in N\)、\(\bar{x}\in\overline{M}\)、\(\bar{x}'\in\overline{M'}\) 满足：(1) \(\bar{f}(\bar{x},\bar{x}')\) 是 y 在 \(\bar{y}\) 中的类 \(\overline{N}\)；(2) \(\overline{N}\) 的每个元素都可写成
\end{enumerate}

\[
\bar {f} (\overline {{{x}}}, \overline {{{z}}} ^ {\prime}) + \bar {f} (\overline {{{z}}}, \overline {{{x}}} ^ {\prime}),
\]

其中 \(\bar{z} \in \overline{\mathbf{M}}\)，\(\bar{z}' \in \overline{\mathbf{M}}'\)。证明若 N 是 Hausdorff 的且 M、M' 完备，则存在 \(x \in \bar{x}\) 和 \(x' \in \bar{x}'\)，使得 \(f(x, x') = y\)（仿照 Hensel 引理的证明作归纳）。在什么条件下 \(x\) 和 \(x'\) 被唯一确定？

\begin{enumerate}
\def\labelenumi{\arabic{enumi}.}
\setcounter{enumi}{1}
\tightlist
\item
  设 A 为局部环，m 为其极大理想，k = A/m 为其剩余域，\(f: A \to k\) 为典范同态。令 \(P \in A[X]\) 为次数为 n 的首一多项式。~记 \(B = A[X]/P \cdot A[X]\)，并以 x 表示 X 在 B 中的典范像。
\end{enumerate}

\begin{enumerate}
\def\labelenumi{(\alph{enumi})}
\item
   设 \(\mathbf{Q}, \mathbf{Q}'\) 是 \(\mathbf{A}[\mathbf{X}]\) 中两个强互素的首一多项式，满足 \(\mathbf{P} = \mathbf{QQ}'\)。证明环 \(\mathbf{B}\) 是理想 \(\mathbf{B}. \mathbf{Q}(x)\) 和 \(\mathbf{B}. \mathbf{Q}'(x)\) 的直和。
\item
  反之，设 \(B = b \oplus b'\) 是 B 作为两个理想直和的分解。证明存在多项式 Q、\(Q'\) 在 A{[}X{]} 中，满足 (a) 的假设，并且 \(\mathfrak{b} = \mathrm{B}.Q(x)\)、\(\mathfrak{b}' = \mathrm{B}.Q'(x)\)。（先证明 b/mb 和 \(b'/mb'\) 由首一多项式 \(Q_{0} \in \bar{f}^{-1}(\mathfrak{b})\)、\(Q'_{0} \in \bar{f}^{-1}(\mathfrak{b}')\) 在 B/mB 中的像生成，并满足
\end{enumerate}

\[
\bar {f} (\mathrm{P}) = \bar {f} (\mathrm{Q} _ {0}) \bar {f} (\mathrm{Q} _ {0} ^ {\prime}).
\]

令 \(r = \deg(Q_0)\)、\(s = \deg(Q_0')\)。证明 b 是以 \(Q_0(x)\)、\(xQ_0(x)\)、\(\ldots\)、\(x^{s-1}Q_0(x)\) 为基的自由 A-模（使用第二章，第 3 节，第 2 节，命题 5）；于是可写成 \(x^sQ_0(x) = a_0Q_0(x) + a_1xQ_0(x) + \cdots + a_{s-1}x^{s-1}Q_0(x)\)，其中 \(a_i \in A\)（\(0 \leqslant i \leqslant s-1\)）；证明若写成

\[
\mathrm{Q} ^ {\prime} (\mathrm{X}) = \mathrm{X} ^ {s} - (a _ {0} + a _ {1} \mathrm{X} + \dots + a _ {s - 1} \mathrm{X} ^ {s - 1}),
\]

则 \(\bar{f}(\mathbf{P}) = \bar{f}(\mathbf{Q}_0)\bar{f}(\mathbf{Q}')\)，且 \(\mathbf{Q}'\in \bar{f}^{-1}(\mathfrak{b}')\)。同样从 \(\mathbf{Q}\) 和 \(\mathbf{Q}'\) 出发定义 \(\mathbf{Q}_0\)，并使用第 1 节命题 2，证明 \(\mathbf{Q}\) 和 \(\mathbf{Q}'\) 解决此问题。）

¶ 3. 设 A 为局部环，m 为其极大理想，k = A/m 为其剩余域，f: A → k 为典范同态。证明以下两个条件等价：

\begin{enumerate}
\def\labelenumi{(\Alph{enumi})}
\setcounter{enumi}{7}
\item[(H)]
  对每个首一多项式 \(\mathbf{P} \in \mathbf{A}[\mathbf{X}]\)，以及 \(\bar{f}(\mathbf{P}) \in k[\mathbf{X}]\) 作为互素首一多项式之积 \(\bar{f}(\mathbf{P}) = \overline{\mathbf{Q}}\)。\(\overline{\mathbf{Q}}'\) 的每个分解，存在 \(\mathbf{Q}\) 中两个首一多项式 \(\mathbf{Q}'\)、\(\mathbf{A}[\mathbf{X}]\)，使得 \(\bar{f}(\mathbf{Q}) = \overline{\mathbf{Q}}\)、\(\bar{f}(\mathbf{Q}') = \overline{\mathbf{Q}}'\)，且 \(\mathbf{P} = \mathbf{QQ'}\)。
\item[(C)]
  每个作为有限生成 A-模的 A 上交换代数，都是局部 A-代数的直积。
\end{enumerate}

（为证明 (H) 蕴含 (C)，仿照第 6 节命题 8，使用习题 2 (a)。为看出 (C) 蕴含 (H)，先证明对于每个作为有限生成 A-模的交换 A-代数 B，B/mB 作为两个理想的直和的每个分解必然形如 b/mb ⊕ b'/mb'，其中 B = b ⊕ b' 是 B 作为两个理想的直和分解；然后使用习题 2 (b)。）

满足条件 (H) 和 (C) 的局部环称为 Hensel 环。每个完备 Hausdorff 局部环都是 Hensel 环。若 A 是 Hensel 环，B 是交换 A-代数、局部环且为有限生成 A-模，则 B 是 Hensel 环。

\begin{enumerate}
\def\labelenumi{\arabic{enumi}.}
\setcounter{enumi}{3}
\tightlist
\item
  \begin{enumerate}
  \def\labelenumii{(\alph{enumii})}
  \tightlist
  \item
    设 \((\mathbf{A}_{\alpha}, \phi_{\alpha \beta})\) 为 Hensel 局部环的直接系统，同态 \(\phi_{\alpha \beta}\) 为局部同态。证明局部环 \(\mathbf{A} = \varinjlim \mathbf{A}_{\alpha}\)（第二章，第 3 节，习题 16）是 Hensel 环（使用习题 3 的判据 (H)）。
  \end{enumerate}
\end{enumerate}

\begin{enumerate}
\def\labelenumi{(\alph{enumi})}
\setcounter{enumi}{1}
\tightlist
\item
  设 K 为交换域，L 为 K 的代数扩张。由 (a) 推出环 \(L \otimes_{K} K[[X_{1}, \ldots, X_{n}]]\) 是 Hensel 环。* (c) 设 A 为 Hensel 局部环，B 为在 A 上整的交换 A-代数（第五章），且 B 是局部环。证明 B 是 Hensel 环（使用 (a)）。*
\end{enumerate}

¶ 5. 设 A 为 Hensel 局部环，B 为有限生成 A-模的（不必交换的）A-代数，b 为 B 的双边理想，令 \(\overline{\mathbf{B}} = \mathbf{B} / \mathfrak{b}\)。

\begin{enumerate}
\def\labelenumi{(\alph{enumi})}
\item
  证明 \(\varepsilon\) 中的每个幂等元 \(\overline{\mathbf{B}}\) 都是 \(\mathbf{B}\) 中某个幂等元的典范像（考虑由单个元素生成的 \(\mathbf{B}\) 的子代数，将其化约到交换情形）。
\item
  令 \((\varepsilon_{n})_{n\geqslant 1}\) 为 \(\overline{\mathbf{B}}\) 中元素的无限序列，使得
\end{enumerate}

\[
\varepsilon_ {i} \varepsilon_ {j} = \delta_ {i j} \varepsilon_ {j}
\]

对于每个指标有序对 \((i,j)\)。证明 B 中存在幂等元的正交序列 \((e_n)_{n\geq 1}\)，使得 \(\varepsilon_{n}\) 都是 \(e_n\) 的典范像，对所有 \(n\)。（仿照代数，第 VIII 章，第 6 节，习题 10 作归纳；注意若 \(e,e^{\prime}\) 是 B 的两个满足 \(ee^{\prime} = 0\) 的幂等元，则 \(e^{\prime} - e^{\prime}e = e''\) 是满足 \(ee'' = e''e = 0\) 的幂等元。）

\begin{enumerate}
\def\labelenumi{(\alph{enumi})}
\setcounter{enumi}{2}
\tightlist
\item
  现在假设 b 是 B 的 Jacobson 根。令 n 为一个整数，\((\varepsilon_{ij})\)（\(1 \leqslant i \leqslant n, 1 \leqslant j \leqslant n\)）为 \(\overline{B}\) 中的元素族，满足 \(\varepsilon_{ij}\varepsilon_{hk} = \delta_{jh}\varepsilon_{ik}\) 且 \(1 = \sum_{i=1}^{n} \varepsilon_{ii}\)。证明 B 中存在元素族 \((e_{ij})\)（\(1 \leqslant i \leqslant n, 1 \leqslant j \leqslant n\)），满足 \(e_{ij}e_{hk} = \delta_{jh}e_{ik}\) 且 \(1 = \sum_{i=1}^{n} e_{ii}\)，并且对于每个指标有序对，\(\varepsilon_{ij}\) 是 \(e_{ij}\) 的典范像。（使用 (b)、代数第 VIII 章第 6 节的习题 11 以及代数第 VIII 章第 1 节的习题 9。）推出若 \(\overline{B}\) 同构于矩阵环 \(\mathbf{M}_{r}(\overline{\mathbf{D}})\)，其中 \(\overline{D}\) 是不必交换的域，则 B 同构于矩阵环 \(\mathbf{M}_{r}(\mathbf{D})\)，其中 D 是有限生成 A-模的 A-代数，且其 Jacobson 根 \(\mathfrak{d}\) 满足 \(D/\mathfrak{d}\) 同构于 \(\overline{D}\)（参见代数，第 VIII 章，第 1 节，习题 9 和 3）。证明若进一步 B 是 A 上的 Azumaya 代数（第二章，第 5 节，习题 14），则 D 也是。~
\end{enumerate}

\begin{enumerate}
\def\labelenumi{\arabic{enumi}.}
\setcounter{enumi}{5}
\tightlist
\item
  给出一个非交换 Artin 环 A 的例子，它带有由双边理想序列 \((a_{n})\) 构成的滤过，满足 \(a_{0} = A\)、\(a_{2} = 0\)，且第 6 节命题 8 对其不成立（参见代数，第 VIII 章，第 2 节，习题 6）。~
\end{enumerate}

¶ 7. (a) 设 A 为交换环，m 为 A 的理想，且 A 关于 m-进拓扑是 Hausdorff 且完备的，\(m/m^{2}\) 是有限生成 A-模。证明 \(A' = A\{X_{1}, \ldots, X_{r}\}\) 上的拓扑是 \(m'\)-进拓扑，其中 \(m' = mA'\)，且 \(m'/m'^{2}\) 是有限生成的 \((A'/m')\)-模（使用第 2 节，第 11 节，命题 14）。特别地，若 A 是 Noether 环，则 \(A'\) 也是 Noether 环。

\begin{enumerate}
\def\labelenumi{(\alph{enumi})}
\setcounter{enumi}{1}
\tightlist
\item
  设 A 为交换 Noether 环，m 为 A 的理想，且 A 关于 m-进拓扑是 Hausdorff 且完备的。令 \(u: A \rightarrow B\) 为从 A 到 Hausdorff 交换拓扑环 B 的连续同态，使 B 成为 A-代数。证明下列条件等价：
\end{enumerate}

(α) B 是 Noether 环，其拓扑是 mB-进拓扑，B 是完备的，且 B/mB 是 A/m 上的有限生成代数。

(β) B 作为拓扑 A-代数同构于 \(\lim B_{n}\)，其中 \((\mathrm{B}_{n})_{n\geqslant1}\) 是离散 A-代数的逆系统，满足当 \(\phi_{nm}:B_{m}\to B_{n}\) 时映射 \(m\geqslant n\) 是满射，\(\phi_{nm}\) 的核为 \(m^{n+1}B_{m}\)，且 \(B_{1}\) 是 A/m 上的有限生成代数。

(\((\gamma)\)) B 作为拓扑 A-代数同构于形如 \(\mathbf{A}\{\mathbf{X}_1,\dots ,\mathbf{X}_r\}\) 的代数对一个闭理想的商。

（为证明 \((\beta)\) 蕴含 \((\gamma)\)，使用第 2 节，第 11 节，命题 14 以及第 2 节，第 8 节，定理 1。）

\begin{enumerate}
\def\labelenumi{\arabic{enumi}.}
\setcounter{enumi}{7}
\tightlist
\item
  设 A 为带线性拓扑的交换环。加法群 \(A[[X_{1},\ldots,X_{p}]]\) 被识别为乘积群 \(A^{N^{p}}\)，并赋予乘积拓扑 T。
\end{enumerate}

\begin{enumerate}
\def\labelenumi{(\alph{enumi})}
\item
  证明 \(\mathcal{T}\) 与 \(\mathbf{A}[[\mathbf{X}_1,\dots ,\mathbf{X}_p]]\) 的环结构相容，并且是使映射 \(f\mapsto \partial f / \partial \mathbf{X}_i\) 连续的线性拓扑。
\item
  对于 \(P = \sum_{e} \alpha_{e,P} X^{e}\) 的每个元素 \(A[[X_{1}, \ldots, X_{p}]]\)（第 1 节，第 6 节的记号）以及每个完备 Hausdorff 拓扑 A-代数 B（第 2 节，习题 28），若元素 \(\mathbf{x} = (x_{1}, \ldots, x_{p}) \in B^{p}\) 满足族 \((\alpha_{e,P}\mathbf{x}^{e})\)（这里对 \(x^{e} = x_{1}^{e_{1}} \ldots x_{p}^{e_{p}}\) 写作 \(e = (e_{1}, \ldots, e_{p})\)）关于 \(N^{p}\) 的有限子集补集滤子在 B 中收敛到 0，则称 x 可代入 P。此族于是可求和，其和记作 \(P(\mathbf{x})\)。证明使 x 可代入 P 的形式幂级数 \(P \in A[[X_{1}, \ldots, X_{p}]]\) 的集合是 \(S_{x}\) 的子环 \(A[[X_{1}, \ldots, X_{p}]]\)，且映射 \(P \mapsto P(\mathbf{x})\) 是从 \(S_{x}\) 到 B 的同态。
\item
  证明若 \(\mathbf{x}\) 可代入 \(\mathbf{P}\)，且 \(\mathbf{y} = (y_1, \ldots, y_p)\) 是 \(\mathbf{B}^p\) 中的元素，并且对 \(y_i\)，\(1 \leqslant i \leqslant p\) 都是拓扑幂零的，则 \(\mathbf{x} + \mathbf{y}\) 可代入 \(\mathbf{P}\)。
\end{enumerate}

特别地，若 B 中存在一个由拓扑幂零元素组成的 0 的邻域，则对于所有 \(P \in A[[X_{1}, \ldots, X_{p}]]\)，可代入 P 的 \(x \in B^{p}\) 的集合 D 是开的，且从 D 到 B 的映射 \(\mathbf{x} \mapsto \mathrm{P}(\mathbf{x})\) 是连续的。

\begin{enumerate}
\def\labelenumi{(\alph{enumi})}
\setcounter{enumi}{3}
\tightlist
\item
  从现在起假设 A 是 Hausdorff 且完备的，并给 \(A[[X_{1},\ldots,X_{p}]]\) 赋予拓扑 T。为了使 \(P_{1},\ldots,P_{q}\) 中的 q 个形式幂级数 \(A[[X_{1},\ldots,X_{p}]]\) 可代入形式幂级数
\end{enumerate}

\[
\mathrm{Q} \in \mathrm{A} [ [ \mathrm{X} _ {1}, \dots , \mathrm{X} _ {q} ] ],
\]

其充要条件是系统 \((\mathbf{P}_{1}(0),\ldots,\mathbf{P}_{q}(0))\) 可代入 Q。

\begin{enumerate}
\def\labelenumi{(\alph{enumi})}
\setcounter{enumi}{4}
\item
  假设 \(\mathbf{x}\) 可代入每个 \(\mathrm{P}_k\)（\(1 \leqslant k \leqslant q\)），且系统 \((\mathrm{P}_1, \ldots, \mathrm{P}_q)\) 可代入 \(\mathbb{Q}\)。证明 \(\mathbf{x}\) 可代入 \(\mathbb{Q}(\mathrm{P}_1, \ldots, \mathrm{P}_q)\)，系统 \((\mathrm{P}_1(\mathbf{x}), \ldots, \mathrm{P}_q(\mathbf{x}))\) 可代入 \(\mathbb{Q}\)，并且当进一步满足下列假设之一时，\(\mathbb{Q}(\mathrm{P}_1(\mathbf{x}), \ldots, \mathrm{P}_q(\mathbf{x})) = ((\mathbb{Q}(\mathrm{P}_1, \ldots, \mathrm{P}_q))(\mathbf{x})\)：( \(\alpha\) ) \(\mathbb{B}\) 中存在理想 \(n\)，使环对 \((\mathbb{B}, n)\) 满足 Hensel 条件，且 \(x_i\)（\(1 \leqslant i \leqslant p\)）属于 \(n\)；( \(\beta\) ) 形式幂级数 \(\mathbb{Q}\) 是受限的。
\item
  取带离散拓扑的 \(\mathbf{B} = \mathbf{A} = \mathbf{F}_2\)，\(\mathbf{P} = \mathbf{X} + \mathbf{X}^2\)、\(\mathbf{Q} = \sum_{n=0}^{\infty} \mathbf{X}^{2^n}\)；则 \(\mathbf{P}\) 可代入 \(\mathbf{Q}\)，\(x = 1\) 可代入 \(\mathbf{P}\) 和 \(\mathbf{Q}(\mathbf{P})\)，\(\mathbf{P}(x)\) 可代入 \(\mathbf{Q}\)，但 \(\mathbf{Q}(\mathbf{P}(x)) = 0\)，而 \((\mathbf{Q}(\mathbf{P}))(x) = 1\)。
\end{enumerate}
% semantic-fragments: legacy-input-seam
\relax{}
